Seluruh tahap \textit{cleaning} dan \textit{preprocessing} berada di
\texttt{src/preprocessing.py}. Untuk mencegah \textit{data leakage}, data dibagi
terlebih dahulu. Setelah itu, setiap statistik (nilai imputasi, batas pencilan,
daftar kategori, rata-rata, dan simpangan baku) dihitung hanya dari
\texttt{X\_train}, lalu diterapkan apa adanya ke \texttt{X\_val} dan
\texttt{X\_test}.

\section{Pembagian Data}

Fungsi \texttt{split} membagi data latih menjadi 80\% \texttt{X\_train} dan 20\%
\texttt{X\_val} secara \textit{stratified} dengan \textit{seed} 42 sehingga
proporsi kelas batal pada kedua bagian sama. \texttt{X\_val} dipakai untuk
evaluasi \textit{hold-out}, sedangkan \textit{k-fold} dijalankan di dalam
\texttt{X\_train}.

\section{\textit{Data Cleaning}}

\subsection{Duplikat}

Baris dengan seluruh fitur identik dihapus dari \texttt{X\_train} melalui
\texttt{drop\_duplicates}. Duplikat membuat sampel tertentu berbobot ganda. Jika
satu salinan masuk ke \textit{fold} latih dan salinan lain masuk ke \textit{fold}
validasi, skor validasi juga menjadi terlalu tinggi. Duplikat pada \texttt{X\_val}
dan \texttt{X\_test} tidak dihapus karena setiap baris pada kedua data itu tetap
harus diprediksi.

\subsection{Nilai Hilang}

Imputasi dilakukan dalam dua tahap. Pada tahap pertama, \texttt{nightly\_rate\_idr}
diisi dengan median harga per pasangan (\texttt{room\_category},
\texttt{arrival\_month}), dan \texttt{days\_before\_arrival} diisi dengan median
per pasangan (\texttt{booking\_channel}, \texttt{arrival\_month}). Harga kamar
bergantung pada tipe kamar dan musim, sedangkan \textit{lead time} bergantung pada
kanal pemesanan. Karena itu, median kelompok diharapkan lebih mendekati nilai
sebenarnya daripada median global. Kedua fitur ini juga merupakan fitur dengan korelasi
terbesar terhadap target.

Pada tahap kedua, nilai yang masih hilang diisi dengan median untuk fitur numerik
dan modus untuk fitur kategorikal. Median dipilih karena tidak terpengaruh
pencilan dan distribusi fitur yang menceng.

Parameter \texttt{zero\_as\_missing} disediakan untuk memperlakukan nilai 0 sebagai
nilai hilang, misalnya harga kamar Rp0. Data latih memuat 413 baris (1,5\%)
dengan \texttt{nightly\_rate\_idr} bernilai Rp0, dan data uji memuat 85 baris.
Harga terendah berikutnya adalah Rp16.000 sehingga Rp0 terpisah jelas dari harga
lainnya. Sebanyak 267 baris Rp0 berasal dari kanal Complimentary, sedangkan 146
baris lainnya berasal dari Online Travel Agent. Dari 301 pemesanan Complimentary,
267 pemesanan berharga Rp0. Rasio pembatalan baris Rp0 hanya 0,7\%, jauh di bawah
rasio pembatalan baris lain (33,4\%).

Pola ini menunjukkan bahwa Rp0 kemungkinan besar menandai kamar gratis, bukan
harga yang tidak tercatat. Jika Rp0 diganti dengan median, informasi tersebut
hilang. SVM RBF ($C=3$, $\gamma=0{,}1$) dengan kedua pilihan menghasilkan skor
yang setara: \textit{macro F1} \textit{hold-out} 0,8384 dan rata-rata 5-\textit{fold}
0,8328 jika Rp0 dipertahankan, serta 0,8376 dan 0,8329 jika Rp0 dianggap nilai
hilang. Selisih keduanya jauh lebih kecil daripada simpangan baku antar-\textit{fold}
(0,006). Karena skornya setara dan Rp0 memiliki makna yang jelas, konfigurasi
akhir mempertahankan Rp0 sebagai nilai asli (\texttt{zero\_as\_missing} kosong).

\subsection{Pencilan}

Nilai \texttt{days\_before\_arrival} dan \texttt{nightly\_rate\_idr} dibatasi
(\textit{capping}) ke rentang IQR
$[Q_1-1{,}5\,\text{IQR},\ Q_3+1{,}5\,\text{IQR}]$ yang dihitung dari
\texttt{X\_train}. Baris dengan nilai ekstrem tidak dibuang karena baris tersebut
tetap valid dan harus diprediksi. Pembatasan ini mencegah nilai ekstrem
mendominasi jarak pada KNN dan SVM RBF serta \textit{gradient} pada Logistic
Regression.

\section{Rekayasa Fitur}

Fungsi \texttt{engineer} membuat fitur baru hanya dari nilai pada baris yang sama.
Karena tidak memakai statistik dari baris lain, fungsi ini aman diterapkan ke data
mana pun.

\begin{table}[H]
	\centering
	\caption{Fitur hasil rekayasa}
	\label{tab:fitur-baru}
	\small
	\begin{tabular}{l p{8.3cm}}
		\toprule
		\textbf{Fitur} & \textbf{Definisi dan alasan} \\
		\midrule
		\texttt{total\_nights} & \texttt{weekend\_nights} + \texttt{weekday\_nights}; lama menginap total \\
		\texttt{total\_guests} & \texttt{num\_adults} + \texttt{num\_children}; jumlah tamu dalam satu pemesanan \\
		\texttt{month\_sin}, \texttt{month\_cos} & \textit{Encoding} siklik bulan kedatangan agar Desember dan Januari berdekatan \\
		\texttt{doy\_sin}, \texttt{doy\_cos} & \textit{Encoding} siklik urutan hari dalam tahun; tanggal yang tidak valid (misalnya 31 Februari) diganti dengan tanggal terakhir pada bulan itu \\
		\texttt{lead\_x\_price} & $\log(1+\texttt{days\_before\_arrival})\times\texttt{nightly\_rate\_idr}/10^6$; interaksi antara dua fitur dengan korelasi terbesar terhadap target \\
		\texttt{date\_rate} & Rasio pembatalan per tanggal kedatangan (\textit{target encoding}, lihat penjelasan di bawah) \\
		\texttt{date\_count} & $\log(1+n_d)$, dengan $n_d$ jumlah pemesanan pada tanggal kedatangan yang sama; perkiraan tingkat keramaian \\
		\bottomrule
	\end{tabular}
\end{table}

Fitur \texttt{date\_rate} merupakan \textit{target encoding} tanggal kedatangan
dengan \textit{smoothing}:
\begin{equation}
	\texttt{date\_rate}(d)=\frac{\sum_{k:\,\text{tanggal}_k=d}y_k+m\,\bar{y}}{n_d+m},
\end{equation}
dengan $n_d$ jumlah pemesanan pada tanggal $d$, $\bar{y}$ rasio pembatalan global,
dan $m=20$. \textit{Smoothing} menarik nilai tanggal yang jarang muncul ke arah
rata-rata global. Fitur ini dihitung dari label sehingga nilainya pada
\texttt{X\_train} dihitung secara \textit{out-of-fold}. \texttt{X\_train} dibagi
menjadi lima \textit{fold}, dan setiap baris mendapat nilai dari tabel yang
dihitung tanpa \textit{fold}-nya sendiri. \texttt{X\_val} dan \texttt{X\_test}
memakai tabel dari seluruh \texttt{X\_train}. Tanpa langkah ini, label setiap
baris ikut masuk ke fiturnya sendiri dan skor validasi menjadi bias.

\lstinputlisting[caption={Rekayasa fitur (\texttt{src/preprocessing.py})},linerange={25-29,31-39}]{../src/preprocessing.py}
\lstinputlisting[caption={\textit{Target encoding} tanggal secara \textit{out-of-fold} (\texttt{src/preprocessing.py})},linerange={84-108}]{../src/preprocessing.py}

\section{\textit{Encoding} dan \textit{Scaling}}

Atribut \texttt{meal\_package} (Room Only $<$ Breakfast $<$ Half Board $<$ Full
Board) dan \texttt{room\_category} (Type A $<$ \dots{} $<$ Type G) diubah dengan
\textit{ordinal encoding} karena kedua atribut memiliki urutan tingkat layanan.
Atribut \texttt{booking\_channel} tidak memiliki urutan sehingga diubah dengan
\textit{one-hot encoding}. Daftar kategorinya dipelajari dari \texttt{X\_train}.
Kolom \texttt{X\_val} dan \texttt{X\_test} kemudian disamakan dengan
\texttt{reindex}, dan kategori yang tidak muncul di \texttt{X\_train} menjadi
vektor nol.

Setelah \textit{encoding}, semua kolom distandardisasi menjadi $(x-\mu)/\sigma$
dengan $\mu$ dan $\sigma$ dari \texttt{X\_train}. KNN dan SVM RBF memakai jarak
Euclidean, sedangkan Logistic Regression dilatih dengan \textit{gradient descent}.
Tanpa standardisasi, harga kamar yang bernilai jutaan akan mengalahkan fitur yang
bernilai satuan. \texttt{Preprocessor} juga menyediakan opsi \texttt{minmax} dan
\texttt{robust}.

\lstinputlisting[caption={\texttt{Preprocessor.fit} dan \texttt{transform} (\texttt{src/preprocessing.py})},linerange={43-48,50-74,76-82}]{../src/preprocessing.py}

\section{Seleksi Fitur dan Reduksi Dimensi}
\label{sec:seleksi}

Kelas \texttt{Pipeline(model, select=k, pca=n)} menambahkan seleksi fitur dan PCA
untuk setiap model. \texttt{Pipeline} di-\textit{fit} ulang pada setiap
\textit{fold} sehingga seleksi fitur dan PCA tidak melihat data validasi.

Seleksi fitur dilakukan dengan bantuan model itu sendiri. Model dilatih sekali
dengan seluruh fitur, lalu fitur diurutkan berdasarkan
\texttt{feature\_importances\_} atau $|\texttt{coef\_}|$. Model kemudian dilatih
ulang dengan $k$ fitur teratas. Untuk SVM, urutan fitur diambil dari $|w|$ SVM
linear ($w=\sum\alpha_i y_i x_i$) karena SVM RBF tidak memiliki bobot eksplisit.

PCA dihitung dari \textit{eigenvector} matriks kovarians. Jika $n<1$, jumlah
komponen dipilih sampai proporsi variansi kumulatif mencapai $n$. Jika $n\ge1$,
PCA memakai tepat $n$ komponen.

\lstinputlisting[caption={PCA dan \texttt{Pipeline} (\texttt{src/preprocessing.py})},linerange={134-151,153-167,169-176}]{../src/preprocessing.py}

\section{Penanganan Ketidakseimbangan Kelas}

Variabel \texttt{METHOD} di \textit{notebook} menentukan penanganan
ketidakseimbangan kelas, yaitu \texttt{None}, \texttt{'smote'}, atau
\texttt{'undersample'}. \textit{Resampling} dijalankan di dalam \texttt{fit}
setelah \textit{preprocessing} dan hanya diterapkan pada data latih sehingga
distribusi kelas pada data validasi tidak berubah. SMOTE menambah sampel sintetis
kelas batal dengan interpolasi antara sampel kelas batal dan tetangga terdekatnya,
sedangkan \textit{random undersampling} membuang sebagian sampel kelas tidak batal
secara acak. Keduanya membuat proporsi kelas pada data latih menjadi 50:50.

Ketiga pilihan dibandingkan dengan SVM RBF ($C=3$, $\gamma=0{,}1$) memakai
prosedur evaluasi pada Bab~4 (Tabel~\ref{tab:imbalance}). Kolom
\textit{precision} dan \textit{recall} mengacu pada kelas batal di
\texttt{X\_val}.

\begin{table}[H]
	\centering
	\caption{Pengaruh penanganan ketidakseimbangan kelas pada SVM RBF}
	\label{tab:imbalance}
	\small
	\setlength{\tabcolsep}{4pt}
	\begin{tabular}{l r c c c c c}
		\toprule
		\texttt{METHOD} & \textbf{Data latih} & \textbf{\textit{Hold-out}} & \textbf{5-\textit{fold} $\mu$} & \textbf{5-\textit{fold} $\sigma$} & \textbf{\textit{Precision}} & \textbf{\textit{Recall}} \\
		\midrule
		\texttt{None}          & 21.857 & \textbf{0,8384} & \textbf{0,8328} & 0,0062 & 0,8090 & 0,7508 \\
		\texttt{'smote'}       & 29.392 & 0,8305 & 0,8258 & 0,0020 & 0,7421 & 0,8166 \\
		\texttt{'undersample'} & 14.322 & 0,8243 & 0,8173 & 0,0036 & 0,7120 & 0,8501 \\
		\bottomrule
	\end{tabular}
\end{table}

Kedua metode \textit{resampling} menaikkan \textit{recall} kelas batal, dari
75,1\% menjadi 81,7\% dengan SMOTE dan 85,0\% dengan \textit{undersampling}.
Namun, \textit{precision} kelas batal turun lebih besar, dari 80,9\% menjadi
74,2\% dan 71,2\%. Model yang dilatih dengan data 50:50 memprediksi lebih banyak
pemesanan sebagai batal: 36,1\% baris \texttt{X\_val} dengan SMOTE dan 39,1\%
dengan \textit{undersampling}, padahal proporsi pembatalan sebenarnya 32,8\%.
Akibatnya, F1 kelas batal hampir tidak berubah (0,7788; 0,7776; dan 0,7749),
sedangkan F1 kelas tidak batal turun dari 0,8979 menjadi 0,8833 dan 0,8737.
\textit{Macro F1} tanpa \textit{resampling} paling tinggi, baik pada
\textit{hold-out} maupun pada 5-\textit{fold}.

Karena itu, konfigurasi akhir memakai \texttt{METHOD = None}. Dengan perbandingan
kelas sekitar 2:1, SVM tanpa \textit{resampling} sudah mendeteksi 75\%
pembatalan. \textit{Resampling} menggeser model ke arah memprediksi batal lebih
sering daripada proporsi sebenarnya. \textit{Undersampling} memang mempercepat pelatihan karena data latihnya
paling sedikit, tetapi skornya paling rendah.
